\chapter{Limitations and Threats to Validity}
\label{sec:limitations}

This thesis should be explicit about what the system and experiments do not establish. The project already has enough complexity that a serious limitations chapter will strengthen the final argument.

\section{Benchmark and Corpus Limitations}

\begin{itemize}
\item The benchmark corpus is authored and curated rather than harvested from unrestricted real-world malware collections.
\item Sample-task pairs cover important reverse-engineering skills, but they cannot span the full diversity of malicious software.
\item Some tasks are intentionally interpretable and may not capture the hardest real-world analyst workflows.
\end{itemize}

\begin{itemize}
\item \TODO{State how representative the experimental corpus is and where it is intentionally simplified.}
\item \TODO{Discuss any bias introduced by using authored samples whose ground truth is known.}
\end{itemize}

\section{Model, Tool, and Infrastructure Variability}

\begin{itemize}
\item Model behavior may drift over time even under the same nominal configuration.
\item Tool outputs can change with version upgrades or environment differences.
\item Dynamic-analysis infrastructure adds host and VM variability that can affect reproducibility.
\end{itemize}

\begin{itemize}
\item \TODO{Document which versions, APIs, and runtime environments were fixed for the final experiments.}
\item \TODO{Explain how bundle reuse and artifact tracking mitigate, but do not eliminate, this problem.}
\end{itemize}

\section{Judging and Measurement Limitations}

\begin{itemize}
\item Judge scores are a proxy for analyst usefulness, not the same thing as human expert review.
\item Some tasks may reward well-structured reporting more than deep mechanistic understanding.
\item Recovered runs and rebuilt experiments improve completeness but may complicate strict independence assumptions.
\end{itemize}

\begin{itemize}
\item \TODO{Discuss how judge quality was validated and where the rubric may still be weak.}
\item \TODO{Explain any remaining mismatch between automated scoring and analyst value.}
\end{itemize}

\section{Scope of Automation Claims}

\begin{itemize}
\item The system assists reverse engineering; it does not eliminate the need for expert analyst review.
\item Tool orchestration can improve evidence gathering without guaranteeing correctness of interpretation.
\item Automatic unpacking or dynamic execution can help analysis but can also alter the artifact under study or introduce new risks.
\end{itemize}

\begin{itemize}
\item \TODO{Be precise about what kinds of analyst effort this system reduces and what it still cannot replace.}
\item \TODO{Address the correctness and security risks of automatically switching to unpacked artifacts.}
\end{itemize}

\section{Ethical and Operational Considerations}

\begin{itemize}
\item Malware-oriented tooling creates dual-use concerns.
\item Automated binary analysis can surface dangerous capabilities or operational details.
\item Safer experimentation may require containment, logging, and restricted dynamic execution.
\end{itemize}

\begin{itemize}
\item \TODO{Add a short ethics discussion covering dual-use risk, containment, and responsible publication choices.}
\item \TODO{State any restrictions on sharing binaries, prompts, or unpacked artifacts in the final appendix or artifact release.}
\end{itemize}
